\chapter{The Direct3D 9/11/12 Backends}
\label{ch:direct3d-backends}

All three of CNA's Windows-only Direct3D backends share one development loop:
cross-compiled from this project's own Debian machine using the same MinGW-w64 toolchain
Chapter~\ref{ch:building-cna}'s Windows cross-build section already introduced, then run and
verified under Wine on a real GPU --- \texttt{D3D9} and \texttt{D3D11} through DXVK,
\texttt{D3D12} through a different translation layer, vkd3d-proton, in its own separate Wine
prefix. None of the three has been verified on real Windows hardware yet; that remains open
for all three alike, tracked as its own \texttt{needs\_human} task in each backend's plan.

\section{D3D9: the one backend that targets pixel authenticity, not just parity}

\texttt{D3D9} is architecturally different in ambition from every other backend in this book.
Its own documentation states the goal directly: ``pixel-for-pixel indistinguishability from
the original XNA~4.0 runtime itself, not just `renders plausibly.''' It runs Microsoft's own
vendored, verbatim XNA~4.0 Stock Effects HLSL, compiled through the real
\texttt{d3dcompiler\_47.dll} rather than a reimplementation of it.

\subsection{The oracle: a real XNA~4.0 runtime as ground truth}

A dedicated tool, \texttt{tools/xna-oracle/}, stands up a genuine, running XNA~4.0 runtime
under Wine --- real GAC assemblies, compiled by the real in-prefix \texttt{csc.exe} --- and
renders the same declarative scene file through both that real runtime and CNA's \texttt{D3D9}
backend, diffing the two resulting images at zero tolerance: an exact match, not an
approximate one. Both sides of the comparison run through the same DXVK translation layer on
the same GPU, which is precisely what makes the comparison a test of CNA's own rendering
correctness rather than of two different graphics stacks. This corpus keeps growing across
this book's own successive passes over this chapter --- 36 scenes as of an earlier pass, 39
real scene files on disk as of this book's own most recent, direct recount --- with zero known
divergence at every count checked so far.

\subsection{What is real}

Full device lifecycle via plain \texttt{Direct3DCreate9} (not D3D9Ex, a deliberate choice),
with genuine \texttt{DeviceLost} / \texttt{DeviceResetting} / \texttt{DeviceReset} handling; all
five XNA stock effects, with 61 of 66 compiled shader variants byte-identical to Microsoft's
own shipped bytecode (the remaining five differ only by HLSL compiler version, and are
separately proven equivalent through the oracle rather than merely assumed so); a real,
oracle- and mutation-verified \texttt{SpriteEffect.fx}, including the half-pixel offset
Direct3D~9 is notorious for --- caught before it could become a false-positive ``closed''
claim, because the offset's absence is invisible on a $1\times1$ test texture and only shows
up at a larger size; and \cnaclass{GraphicsProfile.Reach} / \texttt{.HiDef} enforcement that is
\emph{real} here, consulting an actual \texttt{D3DCAPS9} structure --- the only CNA backend
where that distinction reflects a genuine capability query rather than a decorative property,
as Chapter~\ref{ch:graphicsdevice} already noted. Render targets sampled back as ordinary
textures, and the \texttt{PreferPerPixelLighting}/specular gaps from
Chapter~\ref{ch:stock-effects}'s \cnaclass{GpuDrawParams} discussion, were both real, open
problems earlier in this backend's history and are now fixed --- the specular fix in
particular proved four of five \texttt{PixelLighting} bytecode variants pixel-perfect via the
oracle; the fifth, an untextured vertex-lit variant, remains permanently unreachable, blocked
by the same missing position-only vertex layout that limits the untextured lit bucket on
every backend.

\subsection{What is not, and the one caveat that applies to every number above}

Every \cnaclass{SurfaceFormat} besides \texttt{Color} remains unsupported (a shared,
cross-backend limitation from Chapter~\ref{ch:textures-rendertargets}, not a D3D9-specific
one); \texttt{Texture3D} has no oracle or GPU round-trip coverage at all --- only its
\cnaclass{GraphicsProfile} size-limit behavior is tested; \cnaclass{OcclusionQuery} is not
built for this backend at all, the same gap D3D12 carries; and
\cnaclass{SpriteSortMode.Immediate} / \texttt{.Texture} were confirmed not viable as
deterministic oracle scenes at all --- \texttt{Immediate} is not something a raster diff can
observe, and \texttt{Texture} sorts by an implementation-defined \texttt{GetHashCode()},
which no two implementations are obligated to agree on. The caveat that applies to every
result in this section: \texttt{D3DCAPS9} itself is synthesized by DXVK on this dev loop, not
reported by genuine XNA-era driver hardware --- real-Windows-hardware verification remains
open, as stated above.

\subsection{A worked example: one real oracle scene, quoted in full}

The oracle's own scene format is a small, declarative text file, not a compiled asset or a
piece of C\texttt{++} --- the simplest real scene in the corpus, \texttt{colored3d.scene},
fits in eleven lines and is worth quoting directly rather than only described:

\begin{lstlisting}
width=256
height=256
profile=HiDef
clearcolor=100,149,237,255
vertexcolor=true
lighting=false
primitive=TriangleList
vertex=-0.6,-0.6,0,255,0,0,255
vertex=0.0,0.7,0,0,255,0,255
vertex=0.6,-0.6,0,0,0,255,255
\end{lstlisting}

This declares exactly the scene Chapter~\ref{ch:stock-effects}'s own \cnaclass{BasicEffect}
coverage would predict from the flags alone: \texttt{vertexcolor=true} with
\texttt{lighting=false} is \cnaclass{BasicEffect}'s \texttt{VertexColorEnabled} path with no
lighting applied, drawing a single red/green/blue-cornered triangle (World/View/Projection all
implicitly identity) over a \texttt{CornflowerBlue} clear --- the well-known default XNA clear
color, matching \texttt{Color.CornflowerBlue}'s own real RGBA value exactly
$(100, 149, 237, 255)$. Both the real XNA~4.0 oracle runtime and \texttt{D3D9} render this
identical declarative scene independently, and the zero-tolerance diff checks two things a
screenshot comparison alone would not separate cleanly: the flat-shaded corners prove basic
vertex-color plumbing works at all, while the \emph{center} of the triangle --- genuinely
Gouraud-interpolated between three different corner colors, not any one of them --- proves the
interpolation itself is correct, not merely that three isolated sample points happen to match.
A more elaborate scene in the same corpus, \texttt{fog\_gradient\_quad}
(Chapter~\ref{ch:easygl-backend}'s own negative-\texttt{FogEnd} finding), uses the identical
file format with a handful of additional \texttt{fog*} keys --- the same declarative scheme
scales from this eleven-line triangle up to every one of the corpus's now-39 scenes.

\subsection{Two real bugs the oracle methodology itself found}

Both are worth naming because they demonstrate what a pixel-exact reference actually buys
you over ordinary unit testing. First, a sort-mode Z-clipping bug: \texttt{zFarPlane=1}
mapped a nonzero \cnaclass{SpriteBatch} layer depth outside Direct3D~9's own valid clip-space
Z range, silently clipping away any sprite with a nonzero layer depth entirely --- fixed by
using \texttt{zFarPlane=-1} instead. Second, and more interesting because it is not
\texttt{D3D9}-specific at all: \cnaclass{GraphicsProfile} never actually reached the real
device, because \texttt{Game}'s own \cnaclass{GraphicsDevice} was eagerly default-constructed
(hardcoded to \texttt{Reach}) before \cnaclass{GraphicsDeviceManager} even existed to apply a
game's requested profile --- meaning a game setting \texttt{GraphicsProfile = HiDef;
ApplyChanges();} had no path to the live device at all. This is a shared, cross-backend bug,
only \emph{discovered} while building \texttt{D3D9}'s own tests, because \texttt{D3D9} is the
one backend where the profile actually matters enough to notice.

\subsection{A worked example: why \texttt{zFarPlane} sign alone hid every nonzero-depth sprite}

The Z-clipping bug above is worth working through numerically, since ``\texttt{zFarPlane=1}
instead of \texttt{-1}'' sounds like a cosmetic sign flip until the actual clip-space
consequence is spelled out. \cnaclass{SpriteBatch}'s orthographic projection maps a
\texttt{layerDepth} value in $[0, 1]$ (XNA's own documented range: 0 = front, 1 = back) onto
Direct3D~9's own native clip-space Z range, which is $[0, 1]$ --- \emph{not} the
$[-1, 1]$ range OpenGL-family APIs use. Constructing that projection with
\texttt{zFarPlane=1} produces a near/far mapping where any \texttt{layerDepth} strictly
greater than what the incorrect matrix treats as its own valid far bound gets clipped by the
GPU's own fixed-function Z-clip test --- silently, with no error, exactly like a sprite drawn
behind the camera would be. A sprite at \texttt{layerDepth=0} (the common case, since most
2D games never set it explicitly) was never affected, which is precisely why this bug could
ship unnoticed: it only manifests for a game that actually uses \cnaclass{SpriteSortMode}'s
depth-based ordering with a nonzero layer value, a real but less common \cnaclass{SpriteBatch}
usage pattern than the default. The fix, \texttt{zFarPlane=-1}, corrects the matrix so the
full documented $[0,1]$ \texttt{layerDepth} range maps inside Direct3D~9's own valid clip
volume instead of partly outside it.

\section{D3D11: the first native, dependency-free Windows path}

\texttt{D3D11} is CNA's first backend built directly on native Direct3D~11, verified in the
same MinGW-w64-plus-Wine-plus-DXVK loop as \texttt{D3D9}. It proves out a real
\texttt{ID3D11Device} / \texttt{ID3D11DeviceContext} pipeline: buffers, textures, and render
targets (including MSAA and multiple render targets), state objects, all ten stock HLSL
shader variants across the five stock effects, a real \cnaclass{SpriteBatch}, and --- unique
to this backend and \texttt{D3D12} --- a runtime-\texttt{D3DCompile()}-backed
\texttt{ShaderEffect} path for custom shaders, already introduced in
Chapter~\ref{ch:shader-fx-gap}. All of this is pixel-verified via real GPU readback, not
merely ``returned \texttt{S\_OK}.''

Device-lost/removed recovery here is honestly scoped as \emph{detection-only}: the
device-removed check is real and correctly wired, but has genuinely never fired in this dev
loop, since no real device removal has occurred under Wine+DXVK testing --- full recovery
remains unverified, not because the code is wrong, but because the triggering condition has
never actually happened. One real, structural bug is worth recounting for what it reveals
about static linking: an early build hit an \texttt{undefined reference to Effect::Apply()}
link failure, traced to a genuine circular dependency (the D3D11 \cnaclass{SpriteBatch}
backend calls back into \texttt{Effect::Apply()}, which lives in the main CNA library, which
itself depends on the backend). MinGW's single-pass archive resolution never revisited the
already-searched \texttt{libCNA.a} to satisfy it. The fix was an explicit repeated
\texttt{target\_link\_libraries(\ldots{} CNA)} for the D3D11/D3D12 targets, deliberately
exploiting CMake's own documented support for a static-library link cycle --- and the bug was
confirmed genuinely pre-existing (via \texttt{git stash}, it reproduced identically on the
base commit), not something the D3D11 work itself introduced.

\section{D3D12: built on D3D11's own foundation, and an honest split verification story}

\texttt{D3D12} deliberately reuses \texttt{D3D11}'s own HLSL/DXBC bytecode, format and state
mapping tables, and constant-buffer layouts --- \texttt{D3DCommon} from
Chapter~\ref{ch:backend-architecture} is exactly this shared foundation. It proves out a real
device, command queues, descriptor heaps, command lists, fences with genuine two-frame
back-pressure, a real per-resource barrier-transition tracker, pipeline state objects with
runtime-settable state, root signatures, real per-slot dynamic sampler state, buffers, 2D/cube/%
3D textures, real render targets (including multiple render targets with real mip-chain
generation and device-queried MSAA), real occlusion queries, all ten stock shader variants
(reusing \texttt{D3D11}'s own compiled bytecode directly, not re-deriving it), a real
\cnaclass{SpriteBatch} with the same runtime-compiled custom-shader support as
\texttt{D3D11}, and a real, functionally-proven device-removed recovery path.

The single most important structural fact about this backend's own verification is worth
stating plainly, because it qualifies every number above: \textbf{the routine CTest suite for
this backend is entirely off-screen}. Real swap-chain creation and \texttt{Present()} through
a live window are proven, but only through a separate, manual diagnostic tool run through a
Proton-managed launch --- never through the routine, automated CTest run, because a full
Proton bootstrap was judged too heavy and slow for a normal test run in this dev loop.

\subsection{A genuine architecture mismatch, not a CNA bug}

Swap-chain creation under \emph{plain} Wine (as opposed to a Proton-managed launch) crashes
outright, reproduced twice with a full symbolized backtrace pointing to a null-pointer read
inside Wine's own \texttt{dxgi.dll}. The root cause, once found, is genuinely external to
CNA: a mismatch between Debian's own system \texttt{dxgi.dll} and vkd3d-proton's separately
overridden \texttt{d3d12.dll} --- two DLLs that were never meant to be paired this way. Once
launched through Proton properly, which supplies the matched DLL pair vkd3d-proton actually
expects, swap-chain creation and a genuine ten-frame clear-and-present loop through a real
window both succeed cleanly, verified twice with zero crashes.

\subsection{Explicitly corrected claims}

\texttt{D3D12}'s own documentation is unusually direct about its own earlier revisions being
stale: runtime-settable \cnaclass{BlendState} / \cnaclass{DepthStencilState}/%
\cnaclass{RasterizerState} feeding real pipeline-state objects, sixteen-slot dynamic sampler
state, real public render-target and render-target-cube backends (including multiple render
targets), real \cnaclass{Texture3D}, and real \cnaclass{TextureCube::GetData()} were all
previously described as open gaps in earlier documentation passes and are now all real and
closed --- with the documentation's own closing line telling a reader to re-verify any
\emph{other} specific claim against the project's own live tracking file before trusting it,
rather than this book's snapshot alone.

\section{A reopened feature, and a real architectural divergence hiding under a shared foundation}

\texttt{D3DCommon}'s shared HLSL/DXBC and state-mapping tables, named above, make it easy to
assume the two backends are near-identical beneath that layer. A late pair of tasks in the
Direct3D~11/12 implementation plan --- opened well after both backends were otherwise
considered feature-complete, once a full feature-matrix percentage audit against \texttt{EasyGL}
turned up a handful of remaining gaps --- shows a case where that assumption holds
(\cnaclass{RenderTargetCube} MSAA) and one where it genuinely does not (mip-chain generation).

\subsection{\cnaclass{RenderTargetCube} MSAA: a deliberate exclusion, reopened on request}

Multisampled render-target cubes were, at one point, a matched, deliberate scope exclusion on
both backends, not a silent oversight --- the project owner explicitly asked for it to be
reopened alongside the rest of a late feature-matrix sweep. The two implementations that
resulted are a real worked example of the same underlying Direct3D constraint solved twice,
independently, once per API generation. Neither \texttt{D3D11\_RESOURCE\_MISC\_TEXTURECUBE} nor
a \texttt{D3D11\_SRV\_DIMENSION\_TEXTURECUBE} view can ever be multisampled, and \texttt{D3D12}'s
equivalent \texttt{D3D12\_SRV\_DIMENSION\_TEXTURECUBE} has no multisampled variant either --- so
on both backends, the MSAA color resource becomes a plain, non-cube six-slice
\texttt{Texture2DMSArray} used only as a render-target view, while a second, single-sample,
genuinely cube-flagged resource is resolved into via \texttt{ResolveSubresource()} on unbind, and
it is that resolved resource, not the MSAA one, that the real shader-resource view always
targets. \texttt{D3D12}'s own version of that same resolve additionally requires two explicit
resource-barrier transitions (\texttt{RESOLVE\_SOURCE} on the color resource,
\texttt{RESOLVE\_DEST} on the resolve target) that \texttt{D3D11} does not need at all --- its
context-level \texttt{ResolveSubresource()} call handles the equivalent state transition
implicitly. Both backends resolve only the currently active face on unbind, the same
one-face-at-a-time convention their mip-chain handling (below) also follows. Real proof on both
backends: an $8\times8$ cube face requested at 4x MSAA, cleared, resolved on unbind, and read
back exactly matching, with an explicit check that \texttt{GetMultiSampleCount()} reports the
real requested sample count rather than a silent single-sample fallback.

\subsection{Mip-chain generation: one native GPU call versus an explicit CPU round trip}

Generating a \cnaclass{RenderTargetCube}'s mip chain after a face is drawn to is where the two
backends' shared foundation runs out. \texttt{D3D11} calls
\texttt{ID3D11DeviceContext::GenerateMips()} directly on the cube's own shader-resource view ---
a single native, hardware-accelerated call with no per-face argument at all, since the view it
operates on spans the whole six-face resource. Calling it after any one face is drawn
regenerates the entire cube's mip chain from each face's own current level-0 content in one GPU
call; this is harmless for the five faces that were not just drawn to (their level-0 content did
not change, so recomputing their mips from it is a no-op in effect) and correct for the one that
was. \texttt{D3D12} has no equivalent convenience API for generating mips on an arbitrary format
at all, so \texttt{D3D12RenderTargetCubeBackend::GenerateMipsEXT()} implements the whole chain by
hand: for each mip level, it reads the previous level's pixels back from the GPU for the one
active face only (\texttt{ReadbackSubresourceRGBA8()}), box-filter-downsamples them on the CPU
(\texttt{BoxFilterDownsample()}), and uploads the result into the next level
(\texttt{UploadSubresourceRGBA8()}) --- an explicit, per-level CPU round trip, deliberately scoped
to only the face that was actually just drawn, since repeating that CPU work for all six faces on
every unbind would be needless cost for five faces whose content had not changed. The two
mechanisms reach the same visible result --- a correctly downsampled mip chain for the face that
was drawn --- through genuinely different means: one hardware call operating on the whole
resource at once, and one explicit, face-scoped software loop working level by level. It is a
concrete illustration of what ``\texttt{D3DCommon} is the shared foundation'' does and does not
promise: shared bytecode and state-mapping tables, not identical implementation strategy for
every individual feature, and a reader porting code that depends on either backend's own timing
or cost characteristics for this specific operation should not assume the other backend behaves
the same way underneath.
